
Binary pass/fail reward is the standard signal used in reinforcement learning with execution feedback (RLEF) for code LLMs. CodeRL \citep{Le2022CodeRL}, PPOCoder \citep{Shojaee2023PPOCoder}, RLEF/Gehring et al. [2024], and DAPO \citep{Yu2025DAPO} all assign reward 1 only when all test cases pass and 0 otherwise. Under Group Relative Policy Optimization (GRPO) \citep{Shao2024DeepSeekMath}, this formulation produces zero gradient contribution in any group where all completions fail every test case, because GRPO computes advantages by subtracting the group mean from each completion's reward: when all rewards are 0, all advantages are 0, and no gradient flows.

This dead zone is not an edge case. Under realistic early-training conditions --- a partial pass probability of 0.1 per test case, 8 completions per group, 5 test cases per problem --- 987 of 1,000 simulated groups (98.7\%) satisfy the zero-gradient condition under binary reward. Ratio reward (k/n test cases passing) escapes this regime: it provides non-zero advantage variance whenever at least one completion achieves partial credit and no completion achieves full credit. For the illustrative group with pass counts [0, 0, 1, 2, 0, 3, 0, 0], ratio reward produces advantages spanning [$-0.150$, +0.450] with variance 0.0475, while binary reward produces all-zero advantages with variance 0.0000. This difference is a mathematical consequence of the reward functions, not a probabilistic claim.

DAPO \citep{Yu2025DAPO} acknowledges sparse rewards in GRPO-based code training and introduces mitigations such as dynamic sampling temperature, but retains binary reward and does not quantify the fraction of training groups that produce zero gradient. Prior RLEF work using PPO [Le et al., 2022; Shojaee et al., 2023] is partially insulated from this problem because PPO's value function provides a non-zero baseline even when execution rewards are zero; GRPO eliminates the critic, making the dead zone more consequential.

This paper addresses the gap through a two-phase experimental design. Phase 1 (h-e1) proves the mechanistic claim: ratio reward provides a mathematical guarantee of non-zero advantage variance in groups where binary reward guarantees zero variance, confirmed via synthetic group computation, a 1,000-group simulation, and a 3-step smoke test (exit code 0). Phase 2 (h-m1) tests whether this mechanistic difference translates to policy-level performance differences in a 1,000-step GRPO training run. The h-m1 experiment produced a precision null result: both conditions showed reward\_mean = 0.0 at all 208 logged training steps, with clipped\_ratio = 1.0 and fraction\_partial = NaN throughout --- confirming that 512-token generation is insufficient for DeepSeek-Coder-6.7B-instruct to produce syntactically complete solutions for APPS competition problems, making binary and ratio rewards mathematically equivalent (both zero for all completions at all steps).

The contributions of this paper are as follows. First, a formal proof that ratio reward guarantees non-zero GRPO advantage variance in groups where binary reward guarantees zero variance, with quantification that 98.7\% of early-training groups satisfy this condition under realistic partial-pass distributions (advantage variance: 0.0475 ratio vs. 0.0 binary). Second, the first quantitative analysis of binary reward's gradient dead zone in GRPO training for code LLMs, identifying the exact group-level condition under which the dead zone occurs. Third, characterization of the experimental prerequisite --- non-zero training-set partial solve rate --- required for policy-level ratio-versus-binary differences to be observable, demonstrated via a precision null result with direct diagnostic evidence (clipped\_ratio = 1.0, fraction\_partial = NaN at all 208 logged steps). Fourth, a complete, reproducible GRPO training framework (35/35 unit tests passing across two hypothesis implementations; 9 engineering issues documented and resolved) that enables future experiments under correctly-powered conditions.

